\documentclass[10pt,a4paper]{amsart}
\usepackage{amsmath,amssymb,amsthm,graphicx}
\usepackage{verbatim}
\newenvironment{reference}{\comment}{\endcomment}
%\newenvironment{reference}{}{}
\newenvironment{slogan}{\comment}{\endcomment}
\newenvironment{history}{\comment}{\endcomment}

% For commutative diagrams we use Xy-pic
\usepackage[all]{xy}

% We use 2cell for 2-commutative diagrams.
\xyoption{2cell}
% Theorem environments.
%
\theoremstyle{plain}
\newtheorem{theorem}[subsection]{Theorem}
\newtheorem{proposition}[subsection]{Proposition}
\newtheorem{lemma}[subsection]{Lemma}

\theoremstyle{definition}
\newtheorem{definition}[subsection]{Definition}
\newtheorem{example}[subsection]{Example}
\newtheorem{exercise}[subsection]{Exercise}
\newtheorem{situation}[subsection]{Situation}

\theoremstyle{remark}
\newtheorem{remark}[subsection]{Remark}
\newtheorem{remarks}[subsection]{Remarks}

\numberwithin{equation}{subsection}

% Macros
%
\def\lim{\mathop{\mathrm{lim}}\nolimits}
\def\colim{\mathop{\mathrm{colim}}\nolimits}
\def\Spec{\mathop{\mathrm{Spec}}}
\def\Hom{\mathop{\mathrm{Hom}}\nolimits}
\def\Ext{\mathop{\mathrm{Ext}}\nolimits}
\def\SheafHom{\mathop{\mathcal{H}\!\mathit{om}}\nolimits}
\def\SheafExt{\mathop{\mathcal{E}\!\mathit{xt}}\nolimits}
\def\Sch{\mathit{Sch}}
\def\Mor{\mathop{\mathrm{Mor}}\nolimits}
\def\Ob{\mathop{\mathrm{Ob}}\nolimits}
\def\Sh{\mathop{\mathit{Sh}}\nolimits}
\def\NL{\mathop{N\!L}\nolimits}
\def\CH{\mathop{\mathrm{CH}}\nolimits}
\def\proetale{{pro\text{-}\acute{e}tale}}
\def\etale{{\acute{e}tale}}
\def\QCoh{\mathit{QCoh}}
\def\Ker{\mathop{\mathrm{Ker}}}
\def\Im{\mathop{\mathrm{Im}}}
\def\Coker{\mathop{\mathrm{Coker}}}
\def\Coim{\mathop{\mathrm{Coim}}}

% Boxtimes
%
\DeclareMathSymbol{\boxtimes}{\mathbin}{AMSa}{"02}

%
% Macros for moduli stacks/spaces
%
\def\QCohstack{\mathcal{QC}\!\mathit{oh}}
\def\Cohstack{\mathcal{C}\!\mathit{oh}}
\def\Spacesstack{\mathcal{S}\!\mathit{paces}}
\def\Quotfunctor{\mathrm{Quot}}
\def\Hilbfunctor{\mathrm{Hilb}}
\def\Curvesstack{\mathcal{C}\!\mathit{urves}}
\def\Polarizedstack{\mathcal{P}\!\mathit{olarized}}
\def\Complexesstack{\mathcal{C}\!\mathit{omplexes}}
% \Pic is the operator that assigns to X its picard group, usage \Pic(X)
% \Picardstack_{X/B} denotes the Picard stack of X over B
% \Picardfunctor_{X/B} denotes the Picard functor of X over B
\def\Pic{\mathop{\mathrm{Pic}}\nolimits}
\def\Picardstack{\mathcal{P}\!\mathit{ic}}
\def\Picardfunctor{\mathrm{Pic}}
\def\Deformationcategory{\mathcal{D}\!\mathit{ef}}
\usepackage{cleveref}
\crefname{originalreference}{original reference}{original references}
\Crefname{originalreference}{Original reference}{Original references}
\makeatletter
\newlabel{more-algebra-lemma-fitting-ideal-basics}{{07ZA}{source}{}{sourcepointer-1}{}}
\newlabel{more-algebra-lemma-fitting-ideal-basics@cref}{{[originalreference][0][]07ZA}{[1][0][]source}{}{}{}}
\makeatother
\crefname{originalreference}{original reference}{original references}
\Crefname{originalreference}{Original reference}{Original references}
\makeatletter
\newlabel{more-algebra-lemma-exact-length-1}{{00MX}{source}{}{sourcepointer-2}{}}
\newlabel{more-algebra-lemma-exact-length-1@cref}{{[originalreference][0][]00MX}{[1][0][]source}{}{}{}}
\makeatother
\title{Affine case: \$\textbackslash{}mathrm\{Supp\}(E) \textbackslash{}subset \textbackslash{}mathrm\{Supp\}(F)\$ with F finitely presented forces a nonzero map F → E}
\author{The Stacks Project Authors}
\date{stacks-project:functors.tex:local-snapshot-20261012; GNU FDL 1.2 or later}
\begin{document}
\maketitle
\noindent\textit{Editorial note.} All delivered mathematical text is the author's own, copied verbatim from the pinned Stacks Project chapter file functors.tex; no statement or proof step was rewritten, reordered or removed.\par\medskip
\noindent\textit{Editorial note.} Span recomputation by label pairing (LANE\_BRIEF 10.2): in this chapter the author's proof environment follows the statement environment as a standalone block; the spec covers the complete environment pair (statement span, proof span) as recorded in the pool candidate.\par\medskip
\noindent\textit{Editorial note.} Occupation scan: combined claim\_key + statement\_anchor + generic-line scan over all registration files (active and withdrawn) found 0 hits for functors.tex (fu-occupied.json, 2026-10-14); consistent with the duty manager's six-check dispatch.\par\medskip
\ifdefined\hypertarget\hypertarget{sourcepointer-1}{}\fi
\noindent\textit{Source reference.} Original reference 07ZA: sources/stacks-chow/tags.txt, source lines 10140--10140. Cross-chapter standard prerequisite cited by the author inside the counted span; it is used as a pointer and its own proof is not part of this item.\par\medskip
\ifdefined\hypertarget\hypertarget{sourcepointer-2}{}\fi
\noindent\textit{Source reference.} Original reference 00MX: sources/stacks-chow/tags.txt, source lines 839--839. Cross-chapter standard prerequisite cited by the author inside the counted span; it is used as a pointer and its own proof is not part of this item.\par\medskip
\begin{lemma}
\label{lemma-supported-on-support-pre}
Let $X$ be an affine scheme. Let $\mathcal{F}$ be a finitely presented
$\mathcal{O}_X$-module. Let $\mathcal{E}$ be a nonzero quasi-coherent
$\mathcal{O}_X$-module. If
$\text{Supp}(\mathcal{E}) \subset \text{Supp}(\mathcal{F})$,
then there exists a nonzero map $\mathcal{F} \to \mathcal{E}$.
\end{lemma}

\begin{proof}
Let us translate the statement into algebra. Let $A$ be a ring. Let $M$ be a
finitely presented $A$-module. Let $N$ be a nonzero $A$-module. Assume
$\text{Supp}(N) \subset \text{Supp}(M)$. To show: $\Hom_A(M, N)$ is nonzero.
We may assume $N = A/I$ is cyclic (replace $N$ by any nonzero
cyclic submodule). Choose a presentation
$$
A^{\oplus m} \xrightarrow{T} A^{\oplus n} \to M \to 0
$$
Recall that $\text{Supp}(M)$ is cut out by $\text{Fit}_0(M)$ which
is the ideal generated by the $n \times n$ minors of the matrix $T$.
See More on Algebra, Lemma \ref{more-algebra-lemma-fitting-ideal-basics}.
The assumption $\text{Supp}(N) \subset \text{Supp}(M)$ now means that
the elements of $\text{Fit}_0(M)$ are nilpotent in $A/I$.
Consider the exact sequence
$$
0 \to \Hom_A(M, A/I) \to (A/I)^{\oplus n} \xrightarrow{T^t} (A/I)^{\oplus m}
$$
We have to show that $T^t$ cannot be injective; we urge the reader to
find their own proof of this using the nilpotency of elements of
$\text{Fit}_0(M)$ in $A/I$. Here is our proof.
Since $\text{Fit}_0(M)$ is finitely generated, the nilpotency means that
the annihilator $J \subset A/I$ of $\text{Fit}_0(M)$ in $A/I$ is nonzero.
To show the non-injectivity of $T^t$ we may localize at a prime.
Choosing a suitable prime we may assume $A$ is local and $J$ is still
nonzero. Then $T^t$ has a nonzero kernel by
More on Algebra, Lemma \ref{more-algebra-lemma-exact-length-1}.
\end{proof}

\end{document}
